\documentclass[11pt]{article}

\usepackage[margin=1in]{geometry}
\usepackage{amsmath,amssymb,amsthm}
\usepackage{booktabs}
\usepackage{enumitem}
\usepackage{microtype}
\usepackage[colorlinks=true,allcolors=blue]{hyperref}

\newtheorem{theorem}{Theorem}
\newtheorem{proposition}{Proposition}
\newtheorem{corollary}{Corollary}
\newtheorem{lemma}{Lemma}
\newcommand{\E}{\mathbb{E}}
\newcommand{\Q}{\mathbb{Q}}

\renewcommand{\thesection}{S\arabic{section}}
\renewcommand{\theequation}{S\arabic{equation}}
\renewcommand{\thetheorem}{S\arabic{theorem}}
\renewcommand{\theproposition}{S\arabic{proposition}}
\renewcommand{\thecorollary}{S\arabic{corollary}}
\renewcommand{\thetable}{S\arabic{table}}
\title{Supporting material for\\Myopic Sampling Can Be Maximally Suboptimal in\\Risk-Limiting Financial Audits}
\author{}
\date{}
\begin{document}
\maketitle
This supplement collects the support correction, implementation specialization,
conditional interval-policy results, and numerical material accompanying the
short note. The note's sharp factor-$N$ theorem and cost corollary have
self-contained proofs and do not depend on this material.

\section{Notation and the sharp family}
Let $\pi_i>0$ sum to one, $f_i\in[0,1]$, and $m^\star=\sum_i\pi_i f_i$.
With $R_t$ the remaining set and $A_{t-1}$ the observed contribution, use
strictly positive sampling probabilities and the wealth update
\begin{equation}\label{supp:2}
Z_t=\frac{\pi_{I_t}f_{I_t}}{q_t(I_t)},\qquad
W_t(m)=W_{t-1}(m)\{1+\lambda_t(m)[Z_t-(m-A_{t-1})]\},\quad W_0(m)=1.
\end{equation}
Stakes are predictable, valid (the wealth factors are nonnegative), and
uniformly bounded by $L$. The confidence set is the running intersection of
$\{m\in[0,1]:W_t(m)<1/\delta\}$ and the logical intervals
\begin{equation}\label{supp:3}
J_t=\left[A_t,A_t+\sum_{i\in R_{t+1}}\pi_i\right].
\end{equation}
The stopping time is the first draw with diameter at most $\varepsilon$;
$\operatorname{diam}\varnothing=0$. The oracle is
\begin{equation}\label{supp:1}
q_t^{\rm or}(i)=\frac{\pi_i f_i}{\sum_{j\in R_t}\pi_j f_j}.
\end{equation}
The note's construction chooses rational $\varepsilon>0$ with
\begin{equation}\label{supp:8}
\varepsilon<1/3,\qquad(1+2L\varepsilon)^N<1/\delta,
\end{equation}
and rational $\rho\in(0,1]$, and sets
\begin{equation}\label{supp:9}
a=\varepsilon/2,\qquad w=a/(N-1),
\end{equation}
\begin{equation}\label{supp:10}
\pi_0=1-a,\quad f_0=\frac{\rho w}{1-a};\qquad
\pi_i=w,\quad f_i=1\quad(1\le i<N).
\end{equation}
The two candidates used in the proof are
\begin{equation}\label{supp:11}
m_-=a,\qquad m_+=5a.
\end{equation}
The main theorem gives
\begin{equation}\label{supp:14}
V^\star=1,\qquad \E\tau_{\rm or}=1+\frac{N-1}{1+\rho},
\end{equation}
where $V^\star$ is an infimum over fully supported policies for the fixed
betting rule. Giving item~0 cost one and each other item cost $\kappa\ge1$
gives
\begin{equation}\label{supp:16}
V_c^\star=1,\qquad \E C_{\rm or}=1+\frac{(N-1)\kappa}{1+\rho}.
\end{equation}
Neither value statement asserts general nonattainment.

\section{Strict support does not imply nonattainment}
An earlier draft incorrectly described both infima in \eqref{supp:14} and \eqref{supp:16} as always
unattained under strict support. The sharp bounds require only the values of
these infima. The following example shows why nonattainment must not be
claimed for the full bounded-betting class.

Take $N=2$, $\delta=1/20$, $L=1$, $\varepsilon=1/40$, and
\[
\pi=(79/80,1/80),\qquad f=(1/79,1),\qquad
q_1=(1679/1680,1/1680).
\]
These are the sharp-family parameters $\rho=1$, $a=w=1/80$ and satisfy
$(1+2L\varepsilon)^2=(21/20)^2<20$. Use
\[
\lambda_1(m)=\mathbf1\{m>\varepsilon\},\qquad
\lambda_t(m)=0\quad(t\ge2).
\]
This is a predictable rule using the known tolerance. It is valid for every
possible taint vector: $Z\ge0$ and $m\le1$ imply $1+Z-m\ge0$ whenever the
stake is one, while a zero stake gives factor one. Its uniform cap is one.

If the large item is drawn first, the logical interval is
$[1/80,1/40]$. All its candidates have stake zero, so the combined set is
exactly this interval. If the small item is drawn first, $Z_1=21$ and
$J_1=[1/80,1]$. For every $m>\varepsilon$, the wealth is
$22-m\ge21>20$, so that candidate is rejected. All $m\le\varepsilon$
retain wealth one. The combined set is again $[1/80,1/40]$.

Thus the fully supported policy has $\tau=1$ on both branches, and the true
total $m^\star=1/40$ remains covered on both. Unit costs also give attained
expected cost one. The separate zero-first-stake result below retains its
more restrictive nonattainment conclusion.
\section{The exact two-transaction problem}

Assume throughout this section that the first stake is zero. This makes the first-round betting set all of \([0,1]\). After
sampling item \(i\), the combined interval therefore has diameter
\(1-\pi_i\). Define
\[
 A(\pi,\varepsilon)=\{i\in\{1,2\}:1-\pi_i\le\varepsilon\}.
\]

\begin{proposition}[Complete \(N=2\) solution]
For every fully supported first action \(q\),
\begin{equation}\label{supp:6}
\E_q[\tau]=2-\sum_{i\in A(\pi,\varepsilon)}q(i).
\end{equation}
Hence the full-support infimum is \(2\) if
\(\varepsilon<\min\{\pi_1,\pi_2\}\), and \(1\) otherwise.
When \(A\) is a singleton, the infimum is unattained; when \(A\) is empty
or contains both items, every fully supported first action is optimal.
\end{proposition}

\begin{proof}
The first draw stops exactly on \(A\). Otherwise the sole remaining item is
audited at time two and \eqref{supp:3} is a singleton. Thus
\(\E_q[\tau]=q(A)+2\{1-q(A)\}\), and optimizing gives the claim.
\end{proof}

For
\[
 (\pi_1,\pi_2)=\left(\frac34,\frac14\right),\quad
 (f_1,f_2)=\left(\frac13,1\right),\quad
 \varepsilon=\frac13,\quad\delta=\frac1{20},
\]
the full-support infimum, prop-M, and oracle expectations are respectively
\begin{equation}\label{supp:7}
 1<\frac54<\frac32.
\end{equation}
Under strict full support the infimum remains one but is not attained; putting
probability \(1-\eta\) on item 1 gives expectation \(1+\eta\).

\section{ApproxKelly specialization}
Set \(\delta=1/20\), \(L=5/2\), and
\(\varepsilon=1/(20N)\). Then
\[
 (1+2L\varepsilon)^N
 =\left(1+\frac1{4N}\right)^N<e^{1/4}<2<20.
\]
The two witnesses in \eqref{supp:11} also lie at indices 1 and 5 of the allowed uniform
grid with \(40N+1\) points. Thus the released ApproxKelly construction is a
corollary, but neither its zero initialization nor its grid drives the general
theorem.

Prop-M stops no later than the rank of item \(0\), whose expectation is
\begin{equation}\label{supp:17}
 1+\frac{(N-1)w}{\pi_0+w}
 =1+\frac{a}{\pi_0+w}<\frac1{\pi_0},
\end{equation}
which approaches one. If point scores happen to equal the realized taints,
\(S_i=f_i\), prop-MS is identical to \eqref{supp:1} and inherits both sharp lower bounds.
This is a statement about the constructed scores and outcomes, not a claim that
the point model has a certified zero-error guarantee.

The mathematical update uses the exact denominator in \eqref{supp:2}. The released
code adds a floating-point stabilizer; no claim of bit-for-bit numerical
equivalence for arbitrarily small probabilities is made.
\section{Certified AI intervals: an exact minimax policy}

The negative theorems say that point contributions do not encode terminal
uncertainty. Suppose instead that a model, prior audit, or calibrated procedure
supplies simultaneous bounds, clipped to $[0,1]$, with
\begin{equation}\label{supp:18}
 0\le\ell_i\le f_i\le u_i\le1\qquad(1\le i\le N).
\end{equation}
After auditing a set \(B\), the exact image of this box under the total is
\begin{equation}\label{supp:19}
 C_B^{\rm box}=
 \left[
 \sum_{i\in B}\pi_i f_i+\sum_{i\notin B}\pi_i\ell_i,
 \sum_{i\in B}\pi_i f_i+\sum_{i\notin B}\pi_i u_i
 \right].
\end{equation}
Both endpoints are attainable because the uncertainty set is a Cartesian
product. Define the dollar-weighted uncertainty
\begin{equation}\label{supp:20}
 d_i=\pi_i(u_i-\ell_i),\qquad
 D(B)=\sum_{i\notin B}d_i.
\end{equation}
The observation changes the location of \eqref{supp:19}, but not its width \(D(B)\).

This section uses a separate policy class: items may be selected
deterministically, and $\tau_{\rm box}=\min\{0\le t\le N:D(B_t)\le\varepsilon\}$,
where $B_t$ is the set reviewed after $t$ draws. No importance weighting is
used. The strict-support restriction from Section~S1 therefore does not apply
to the box-only certificate.

\begin{theorem}[Pathwise minimax box policy]
Let \(d_{(1)}\ge\cdots\ge d_{(N)}\), and define
\begin{equation}\label{supp:21}
 k^\star=\min\left\{k\in\{0,\ldots,N\}:
 \sum_{r=k+1}^N d_{(r)}\le\varepsilon\right\}.
\end{equation}
Auditing in nonincreasing order of \(d_i\) stops \eqref{supp:19} in exactly
\(k^\star\) reviews. Moreover, over all adaptive randomized policies,
\begin{equation}\label{supp:22}
 \inf_q\sup_{f\in\prod_i[\ell_i,u_i]}
 \E_f[\tau_{\rm box}(q)]=k^\star.
\end{equation}
The lower bound holds on every sample path.
\end{theorem}

\begin{proof}
The sorted policy leaves the residual in \eqref{supp:21}. Any history of length
\(t<k^\star\) can remove at most \(\sum_{r=1}^t d_{(r)}\). Its residual width
therefore exceeds \(\varepsilon\), regardless of observations, randomization,
or adaptation. No policy can stop earlier on any path, while the sorted
construction attains the bound.
\end{proof}

This gives a precise answer to the implementable robust-certificate problem:
audit uncertainty, not merely predicted error. For a multiplicative guarantee,
write the condition without division as
\begin{equation}\label{supp:23}
 (1-b)f_i\le S_i\le(1+b)f_i,\qquad 0\le b<1.
\end{equation}
It implies
\begin{equation}\label{supp:24}
 \ell_i=\frac{S_i}{1+b},\qquad
 u_i=\min\left\{1,\frac{S_i}{1-b}\right\}.
\end{equation}
If no upper endpoint clips, then
\begin{equation}\label{supp:25}
 d_i=\pi_i S_i\frac{2b}{1-b^2}.
\end{equation}
Thus \(\pi_iS_i\) is the correct priority, but it should be used as a
descending deterministic priority for this certificate, not as a randomized
probability. When \eqref{supp:24} clips, the exact priority in \eqref{supp:20} can differ even in
rank from \(\pi_iS_i\).

\begin{corollary}[Point accuracy versus certified audit value]
On the population \eqref{supp:9}--\eqref{supp:10}, let the point scores happen to equal the realized
taints, \(S_i=f_i\), and supply the simultaneous intervals
\begin{equation}\label{supp:26}
 [\ell_0,u_0]=\left[0,\frac{2\varepsilon}{\pi_0}\right],
 \qquad [\ell_i,u_i]=[0,1]\quad(1\le i<N).
\end{equation}
The certified policy stops in one review. Randomized prop-MS has
\[
 \E[\tau_{\rm propMS}]=1+\frac{N-1}{1+\rho}.
\]
With \(c_0=1\) and \(c_i=\kappa\), its expected review cost is \eqref{supp:16}, while the
certified cost optimum is one.
\end{corollary}

\begin{proof}
The large upper endpoint is below one because \(\varepsilon<1/3\). The
intervals contain the realized taints: in particular,
\(\pi_0f_0=\rho w\le2\varepsilon\). Their uncertainty contributions satisfy
\[
 d_0=2\varepsilon,
 \qquad \sum_{i=1}^{N-1}d_i=\sum_{i=1}^{N-1}w
 =a=\frac{\varepsilon}{2}.
\]
No audited set omitting item \(0\) can stop, because its residual contains
\(d_0>\varepsilon\). Auditing item \(0\) leaves residual
\(a\le\varepsilon\), so descending uncertainty stops after one review. For the box stopping rule, \(S=f\) makes prop-MS use the oracle rates,
and it stops exactly at item \(0\)'s rank.
The review and cost formulas follow from the exponential-clock calculation.
\end{proof}

\subsection{Heterogeneous costs}

For general positive \(c_i\), the exact robust problem becomes
\begin{equation}\label{supp:27}
 \begin{aligned}
 \min_{x\in\{0,1\}^N}\quad&\sum_i c_i x_i\\
 \text{s.t.}\quad&\sum_i d_i x_i
 \ge D(\varnothing)-\varepsilon.
 \end{aligned}
\end{equation}
The accompanying verification archive implements an exact rational Pareto-frontier dynamic program.
After each item, a state records removed width and cost, caps width at the
target, and discards a state if another has at least as much removed width at
no greater cost. Induction over items proves that every discarded partial
solution is dominated and that the least-cost terminal state solves \eqref{supp:27}.

\subsection{Risk accounting and a hybrid audit}

Suppose \eqref{supp:18} holds simultaneously with probability at least
\(1-\delta_{\rm AI}\). On that event, \eqref{supp:19} covers \(m^\star\) at every time,
under any adaptive policy. If a betting CS has simultaneous failure probability
\(\delta_{\rm CS}\), their running intersection fails with probability at most
\begin{equation}\label{supp:28}
 \delta_{\rm AI}+\delta_{\rm CS}
\end{equation}
by a union bound; independence is unnecessary. A deployable architecture can
therefore individually examine a certainty stratum chosen by \eqref{supp:20}, then run a
randomized risk-limiting procedure on the residual population, explicitly
allocating the two risk budgets.

\section{Finite-grid Bellman verification}

For general \(N\), the state must contain the ordered history, remaining set,
wealth at every candidate \(m\), ApproxKelly payoff accumulators, and current
running intersection. If \(\Phi(x,q,i)\) is the transition, the exact equation
is
\begin{equation}\label{supp:29}
 V(x)=
 \begin{cases}
 0,&x\text{ is in the stopping region},\\
 1+\displaystyle\inf_q\sum_{i\in R(x)}q(i)V(\Phi(x,q,i)),&\text{otherwise}.
 \end{cases}
\end{equation}
The action changes both transition probability and confidence state. This is
the continuation value absent from the one-step surrogate.

We implement \eqref{supp:29} with exact rational arithmetic on a finite candidate grid
and the strict-support mesh \(q_i\in\{1/D,\ldots,D/D\}\). Backward induction
over the finite horizon gives the global optimum over that finite action set;
no simulation enters these values. Table~\ref{tab:small} uses a 21-point
candidate grid and \(D=6\). Every policy's admissible betting range uses the
unrevealed true taints (\texttt{payoff\_information="oracle"}). These are
fixed-population comparisons with oracle information, not evaluations of
implementable betting policies. The mesh-priority comparator randomizes on
the strict-support mesh according to certified-width priorities; it is not
the deterministic optimal box policy. Policies not lying on this action mesh can beat its
optimum, so the table is verification and structure discovery, not a
continuous-action optimality claim.

\begin{table}[t]
\centering
\small
\caption{Expected stopping times on the finite grid/mesh, rounded from exact
rationals. All policies use oracle payoff bounds. Mesh priority is randomized.}
\label{tab:small}
\begin{tabular}{lrrrrr}
\toprule
Case & Oracle & prop-M & prop-MS & Mesh priority & Mesh Bellman\\
\midrule
Terminal value, \(N=3\) & 2.4167 & 2.1500 & 2.3158 & 2.0556 & 2.0556\\
Mixed information, \(N=4\) & 3.8558 & 3.4488 & 3.8302 & 3.7176 & 3.4213\\
Clipped score, \(N=4\) & 3.2008 & 3.4488 & 3.1839 & 3.4213 & 3.4213\\
\bottomrule
\end{tabular}
\end{table}

The first mesh action in the mixed \(N=4\) case is
\((1/6,1/6,1/2,1/6)\), concentrating on neither the largest recorded value nor
the largest contribution alone. This is numerical evidence for a blended
continuation-value index, not yet a theorem.

\section{Reproducible stress tests}

We generated four deterministic \(N=1000\) synthetic populations from fixed
seeds. Each AI interval follows \eqref{supp:24}; the stopping rule is solely the
simultaneous-box width. The sorted optimum is exact. Comparator means use
1,000 independent exponential-race permutations, which exactly generate their
probability-proportional-without-replacement orders. Table~\ref{tab:synthetic}
reports mean reviews; Monte Carlo standard errors for prop-MS range from 0.17
to 0.30 reviews.

\begin{table}[t]
\centering
\small
\caption{Synthetic certified-box workload; not a client-data claim.}
\label{tab:synthetic}
\begin{tabular}{lrrrrr}
\toprule
Scenario & Sorted optimum & prop-MS & Oracle & prop-M & Uniform\\
\midrule
Unclipped, concentrated & 347 & 452.13 & 452.68 & 493.41 & 895.66\\
Clipped, high risk & 412 & 536.52 & 531.01 & 528.11 & 898.89\\
Diffuse values & 740 & 819.31 & 819.41 & 895.66 & 900.44\\
Tight certificate & 659 & 736.39 & 736.64 & 770.81 & 975.57\\
\bottomrule
\end{tabular}
\end{table}

The exact policy uses 9.7--23.3\% fewer reviews than randomized prop-MS in
these scenarios. These are controlled stress tests of the theorem's mechanism,
not estimates of savings in any audit population. Real validation requires
authorized data, simultaneous calibration checks, review-cost measurement,
and comparison under the applicable audit objective.

\section{Verification artifacts}
The accompanying verification archive contains two exact rational JSON
certificates, independent verification scripts, the complete two-item
characterization, and tests for the arbitrary-$N$ witness inequalities and
cost formulas. The support-attainment example in Section~S2 has a separate
exact rational regression check. Certificates corroborate finite instances;
they do not mechanically certify the general theorem.

The synthetic tables use simulated random sampling orders and are labeled as
such. The finite-grid table instead uses exact rational backward induction
on its explicitly stated oracle-information model. Neither table establishes
field performance or solves the unrestricted continuous-action problem.

\begingroup\small\raggedright
\begin{thebibliography}{2}
\bibitem{shekhar2023}
S.~Shekhar, Z.~Xu, Z.~Lipton, P.~Liang, and A.~Ramdas.
Risk-limiting financial audits via weighted sampling without replacement.
\emph{Proceedings of UAI}, PMLR 216:1932--1941, 2023.
\url{https://proceedings.mlr.press/v216/shekhar23a.html}.
\bibitem{implementation}
S.~Shekhar et al. \emph{WeightedWoRConfSeq}, commit \texttt{a834e459a47f}.
\url{https://github.com/sshekhar17/WeightedWoRConfSeq}.
\end{thebibliography}
\endgroup
\end{document}
